\section{Cùng kiến trúc, đổi dữ liệu}
\label{sec:ch11-hai-nua-cua-bang}

\index{Vision Transformer!ImageNet và JFT}%
Đọc bài ViT chỉ qua con số cao nhất là cách nhanh nhất để bỏ lỡ luận điểm. Bài
có một headline 88.55, nhưng headline ấy đi cùng fine-tune độ phân giải 512 và
Polyak averaging. Nó không phải phép so chỉ đổi dữ liệu.

Phép so sách này cần nằm ở bảng 5, phụ lục C. Cùng ViT-B/16, cùng fine-tune ở
384, đổi tập pretraining từ ImageNet-1k sang JFT-300M:

\begin{bridge}{bảng 5 của Dosovitskiy và cộng sự}
\begin{minted}{text}
pretraining data  ImageNet top-1
ImageNet-1k       77.91
JFT-300M          84.15
\end{minted}

ImageNet-1k có 1,000 lớp và 1.3M ảnh. JFT có 18,000 lớp và 303M ảnh; chữ
\enquote{300M} trong tên là làm tròn. Hai số accuracy là số tác giả báo cáo ở
cùng bảng, cùng kiến trúc, cùng cách fine-tune~\autocite{dosovitskiy2021vit}.
\end{bridge}

Từ 77.91 lên 84.15 không chứng minh một định luật rằng thêm dữ liệu luôn thắng
mọi thiên kiến. Nó nói hẹp hơn và đúng hơn: trong recipe của bài, kiến trúc mà
trước đó không đủ cạnh các baseline tích chập trở thành kiến trúc mạnh khi có
JFT. Chính mục 4.3 của bài mô tả các ViT lớn ban đầu kém hơn Base trên
ImageNet, gần nhau hơn sau ImageNet-21k, rồi bộc lộ lợi thế ở JFT-300M.

Đây cũng là câu trả lời cho thắc mắc hợp lý ở cuối chương trước. Nếu attention
biểu diễn được convolution, tại sao cần convolution? Câu trả lời không phải
là attention thiếu phép tính. Nó thiếu một lời hứa về dữ liệu: ở đây là một
chi tiết gần một pixel đáng được đọc như chi tiết gần pixel khác. JFT cho mô
hình quá nhiều ví dụ để tự mua lại một phần lời hứa ấy.

Nhưng JFT không phải thứ tôi có trên laptop, và một bảng 303M ảnh không biến
thành thí nghiệm bằng cách viết ít layer hơn. Mục sau chỉ đo phần đối nghịch
với headline: điều gì xảy ra khi giữ dữ liệu nhỏ.
